\documentclass[11pt,a4paper]{article}
\usepackage[utf8]{inputenc}
\usepackage{amsmath,amssymb,amsthm}
\usepackage{hyperref}
\usepackage{geometry}
\geometry{margin=1in}

\newtheorem{theorem}{Theorem}
\newtheorem{lemma}[theorem]{Lemma}
\newtheorem{conjecture}[theorem]{Conjecture}
\theoremstyle{definition}
\newtheorem{definition}[theorem]{Definition}
\newtheorem{remark}[theorem]{Remark}

\title{On the Optimality of Varshamov-Tenengolts Codes among Perfect Single-Deletion-Correcting Codes: A Formalization Report}
\author{Austin Anderson \\ \small\href{https://github.com/aanderson3456/MyLean4Code}{GitHub Repository: aanderson3456/MyLean4Code}}
\date{\today}

\begin{document}

\maketitle

\begin{abstract}
This report surveys the single-deletion-correcting properties of binary Varshamov-Tenengolts (VT) codes, focusing on the unsolved conjecture regarding the global optimality of $VT_0(n)$ and the recent Lean 4 formalization efforts. We describe the formal proof that $VT_0(n)$ is optimal among all \emph{perfect} single-deletion-correcting codes, which successfully compiles in the interactive theorem prover Lean 4. While the core maximality of the zero-checksum class is fully verified using complex analysis and discrete Fourier transforms, the bridge showing that any perfect code has size matching a VT code is currently postulated as a conjecture.
\end{abstract}

\section{The VT(0) Optimal Theorem and Conjecture}
For any block length $n \ge 1$ and residue class $a \in \{0, 1, \dots, n\}$, the binary Varshamov-Tenengolts code $VT_a(n)$ consists of all binary vectors $(x_1, \dots, x_n) \in \{0, 1\}^n$ satisfying the checksum condition:
\begin{equation}
\sum_{i=1}^n i x_i \equiv a \pmod{n+1}.
\end{equation}
A survey by N. J. A. Sloane~\cite{Sloane02} outlines the fundamental properties of these codes, showing that $VT_0(n)$ possesses the maximum cardinality among all $VT_a(n)$ codes for a given length $n$. More significantly, it is conjectured (Conjecture~2.6 in~\cite{Sloane02}) that $VT_0(n)$ is optimal among \emph{all} binary single-deletion-correcting codes of length $n$, meaning its cardinality $|VT_0(n)|$ achieves the absolute upper bound $A(n, 1)$. This conjecture has been verified computationally for small lengths: Albert No~\cite{No19} proved the conjecture for $n \le 10$ using Mixed Integer Linear Programming (MILP) techniques, and Nakasho et al.~\cite{Nakasho23} established the bound for $n = 11$ (yielding $|VT_0(11)| = 172$) via an integer linear programming solver with run-based linear constraints. For general $n$, the conjecture remains unsolved.

\section{Lean Formalizations by Yuki Kondo}
The initial formalization of Varshamov-Tenengolts codes in the Lean theorem prover was conducted by Yuki Kondo (c/o Manabu Hagiwara)~\cite{Kondo22}. Kondo formalized the core definitions of binary vectors, the single-deletion deletion sphere $dS(x)$, and the checksum moment function $\sum_{i=1}^n i x_i$ recursively using 1-based indexing. A major milestone of their work was the formalization and verification of Levenshtein's decoding algorithm. Kondo proved that for any codeword $x \in VT_a(n)$ and any received vector $y \in dS(x)$, the decoding algorithm successfully reconstructs the original vector $x$, formally establishing the single-deletion-correcting capability of VT codes in Lean.

\section{Other Attempts and Partial Results}
Various historical and modern attempts have sought to resolve the optimality of VT codes. Levenshtein~\cite{Lev65} established the asymptotic optimality of these codes by showing that $A(n, 1) \sim 2^n / n$ as $n \to \infty$. Furthermore, Levenshtein~\cite{Lev92} introduced the concept of perfect deletion-correcting codes, proving that the deletion spheres $dS(x)$ of any VT code $VT_a(n)$ form a partition of the space of descendants $\{0,1\}^{n-1}$, making VT codes perfect single-deletion-correcting codes for all $a$. Other related efforts, such as those by Tenengolts~\cite{Ten76}, have focused on constructing linear single-deletion-correcting codes, though they suffer from a larger check-bit overhead of $O(\sqrt{n})$ compared to the $O(\log n)$ overhead of VT codes.

\section{Detailed Proof of \texttt{vt0\_is\_optimal\_perfect}}
The theorem \texttt{vt0\_is\_optimal\_perfect}, formalized in the module \texttt{VTlean.Progressive}, establishes that $VT_0(n)$ is optimal among all perfect single-deletion-correcting codes of length $n$.
\begin{theorem}
Let $C \subset \{0,1\}^n$ be a perfect single-deletion-correcting code of length $n$. Then:
\begin{equation}
|C| \le |VT_0(n)|.
\end{equation}
\end{theorem}

\begin{proof}[Proof Structure and Formalization Details]
The proof is structured as follows and builds successfully in Lean 4:
\begin{enumerate}
\item \textbf{Perfect Card Equivalence}: We show that the cardinality of any perfect single-deletion-correcting code $C$ must match the cardinality of some VT code $VT_a(n)$ for some $a \in \{0, \dots, n\}$. Let $c(x) = 1$ if $x \in C$ and $0$ otherwise. The perfect code condition implies that for all descendants $y \in \{0,1\}^{n-1}$, we have $\sum_x A_{x,y} c(x) = 1$, where $A_{x,y} = 1$ if $y \in dS(x)$ and $0$ otherwise. 
By inductive tracking of the run transitions using cumulative deletion potentials $P(x, k, u) = \sum_{y \in \text{cumulativeDels}(x, k)} u(y)$, we unroll this algebraic system. This reduces the size of $C$ to $\sum_{k=0}^n T_{n,k,a}$ for some $a$, which is the size of $VT_a(n)$. In the Lean codebase, this corresponds to the lemma \texttt{perfect\_code\_card\_eq\_sum\_T} in \texttt{VTlean.AlgebraicFractal}, which is currently postulated with a \texttt{sorry}.
\item \textbf{Maximality of the Zero-Checksum Class}: We prove that $|VT_a(n)| \le |VT_0(n)|$ for all $a$. This corresponds to the theorem \texttt{VTCode\_zero\_is\_max}, which is fully verified in the module \texttt{VTlean.Cuculiere} without any occurrences of \texttt{sorry}. The proof utilizes the Discrete Fourier Transform over $\mathbb{C}$. The card of $VT_a(n)$ is represented using the $(n+1)$-th roots of unity $\omega$:
\begin{equation}
|VT_a(n)| = \frac{1}{n+1} \sum_{j=0}^n \omega^{-aj} P_n(\omega^j),
\end{equation}
where $P_n(y) = \prod_{k=1}^n (1 + y^k)$ is the generating polynomial of Cuculiere's triangle. By complex conjugation symmetry, we show that $P_n(\omega^j)$ is a non-negative real number for all $j$. Since $|\omega^{-aj}| = 1$ and $P_n(\omega^j) \ge 0$, the real part of each term in the sum is bounded by $P_n(\omega^j)$. Thus, the sum is maximized when $a = 0$ (yielding $\omega^0 = 1$), concluding that $|VT_a(n)| \le |VT_0(n)|$.
\end{enumerate}
\end{proof}

\begin{thebibliography}{9}
\bibitem{Sloane02}
N. J. A. Sloane, ``On Single-Deletion-Correcting Codes,'' \emph{Codes and Designs} (Ray-Chaudhuri Festschrift), Walter de Gruyter, Berlin, 2002, pp. 273--291.
\bibitem{No19}
A. No, ``Nonasymptotic Upper Bounds on Binary Single Deletion Codes via Mixed Integer Linear Programming,'' \emph{Entropy}, vol. 21, no. 9, p. 868, 2019.
\bibitem{Nakasho23}
K. Nakasho, M. Hagiwara, A. Anderson, and J. B. Nation, ``The Tight Upper Bound for the Size of Single Deletion Error Correcting Codes of Length 11,'' \emph{Proceedings of the 2024 International Symposium on Information Theory and Its Applications (ISITA2024)}, 2024 (also arXiv:2309.14736).
\bibitem{Kondo22}
Y. Kondo (c/o M. Hagiwara), ``Formalization of Varshamov-Tenengolts Codes in Lean,'' GitHub Repository, 2022.
\bibitem{Lev65}
V. I. Levenshtein, ``Binary codes capable of correcting deletions, insertions, and reversals,'' \emph{Soviet Physics Doklady}, vol. 10, no. 8, pp. 707--710, 1966.
\bibitem{Lev92}
V. I. Levenshtein, ``Perfect codes in deletion and insertion metric spaces,'' \emph{Discrete Mathematics and Applications}, vol. 2, no. 3, pp. 241--258, 1992.
\bibitem{Ten76}
G. Tenengolts, ``Codes correcting a single deletion or insertion,'' \emph{IEEE Transactions on Information Theory}, vol. 22, no. 3, pp. 338--345, 1976.
\end{thebibliography}

\end{document}
